\section{Related Work}
\label{sec:related}

\textbf{Robotics VLMs and trajectory-scale data.}
Embodied VLMs and VLAs such as PaLM-E, RT-2, and OpenVLA integrate visual-language representations with robot data to improve planning or action generation~\cite{driess2023palme,brohan2023rt2,kim2024openvla}. These models are enabled by large manipulation datasets such as Open X-Embodiment and DROID, which broaden the diversity of robots, tasks, and scenes available for learning~\cite{oneill2023openx,khazatsky2024droid}. Robo2VLM takes a complementary direction: it uses robot trajectories to generate VQA examples for evaluating and improving VLMs rather than directly training a control policy~\cite{chen2025robo2vlm}.

\textbf{Visual and spatial reasoning benchmarks.}
Classic VQA and compositional reasoning datasets showed that aggregate accuracy can hide question biases and reasoning failures~\cite{johnson2017clevr,hudson2019gqa}. Spatial reasoning benchmarks such as VSR, SpatialVLM, and SpatialRGPT further show that VLMs struggle with relative position, orientation, distance, and 3D-aware region reasoning~\cite{liu2022vsr,chen2024spatialvlm,cheng2024spatialrgpt}. Our benchmark slice follows this diagnostic spirit, but evaluates robotics images and manipulation-specific questions rather than web images or synthetic scenes.

\textbf{Grounding and hallucination checks.}
Multimodal benchmarks such as MME test perception and cognition across broad tasks~\cite{fu2023mme}, while hallucination benchmarks such as POPE probe whether LVLM outputs stay consistent with image content~\cite{li2023pope}. Our sanity probes use the same principle in a robotics VQA setting: after a model answers or is given an answer, we test whether it can reject unsupported claims under blank, shuffled, or deliberately wrong-answer controls.